\chapteropen{07}

\begin{cele}
  \item Zmierzyć własne tempo czytania prostym testem ze stoperem i policzyć z tego wynik w słowach na minutę.
  \item Odróżnić typowe tempo czytania dorosłych od obietnic szybkiego czytania i wiedzieć, kiedy obietnica przekracza to, co da się uczciwie sprawdzić.
  \item Wskazać, co realnie zwiększa liczbę przeczytanych książek: czas, mniej przerwań, dobór lektury i ciche miejsce.
  \item Ocenić, czy subwokalizacja jest problemem, i nie wyrzucać jej w imię mitu.
  \item Czytać trudny tekst wolno, z ołówkiem, i wiedzieć, po czym poznać, że tempo trzeba zmienić.
  \item Policzyć, co daje wzrost tempa, a co dodatkowe dziesięć minut dziennie, na własnych liczbach.
\end{cele}

Tempo czytania jest jedną z dwóch zmiennych, które decydują o tym, ile książek mieści się w roku. Druga to czas, który poświęcasz na czytanie każdego dnia. Rozdział~1 pokazał, że mnożą się one przez siebie: liczba słów na minutę razy liczba minut razy liczba dni daje liczbę słów w roku, a ta dzielona przez długość książki daje liczbę książek. Ten rozdział dotyczy tego pierwszego czynnika, ale z dwoma zastrzeżeniami. Po pierwsze, tempo jest mniej stałe, niż sugerują kursy, i zależy od tekstu. Po drugie, najczęściej nie jest ono wąskim gardłem. Zanim zaczniesz je poprawiać, warto wiedzieć, czy w ogóle jest czego poprawiać.

\section{Zmierz, zanim zaczniesz poprawiać}

Nie znasz swojego tempa, jeśli go nigdy nie mierzyłeś. Większość ludzi szacuje je na podstawie wrażenia (\glqq czytam wolno\grqq{} albo \glqq czytam szybko\grqq{}), a wrażenia w tej sprawie są słabe. Wolne czytanie trudnego tekstu bywa odbierane jako dowód niskiej sprawności, choć jest rozsądną strategią. Szybkie przewracanie kartek bywa odbierane jako sukces, choć po fakcie nic z niego nie zostaje.

Test jest prosty i wymaga trzech rzeczy: książki, kartki papieru, stopera. Wybierz tekst, który czytasz zwykle, czyli zwykłą prozę popularną albo esej, a nie coś, co znasz na pamięć ani nie czytałeś nigdy. Otwórz książkę na początku czytanego fragmentu i zaznacz pięć pełnych stron. Pełnych, czyli takich, na których jest tekst od góry do dołu, bez rozdziałów, ilustracji i stron tytułowych.

Policz słowa. Nie musisz liczyć każdego słowa co do jednego: wystarczy policzyć słowa w pełnym wierszu na jednej stronie, pomnożyć przez liczbę wierszy i uśrednić po pięciu stronach. Dla przykładu przyjmijmy, że na pięciu stronach znalazło się w sumie 1\,480 słów, czyli średnio 296 na stronę. To mniej więcej tyle, ile zwykle mieści się na stronie książki w formacie zbliżonym do tego, w którym czytasz.

Teraz czytanie. Włącz stoper i czytaj tak, jak czytałbyś w zwykły wieczór: bez przyspieszania i bez zwalniania na siłę. Zatrzymaj stoper na ostatnim słowie piątej strony. Załóżmy, że zajęło to 6 minut i 10 sekund, czyli 370 sekund.

Obliczenie jest krótkie:

\begin{przyklad}[Przykład: pięć stron w 370 sekund]
Czas w minutach: $370 / 60 \approx 6{,}17$ min.

Tempo: $1\,480 / 6{,}17 \approx 240$ słów na minutę.

Czyli przy tym tekście i tym czytelniku tempo wynosi 240~sł./min. Liczba jest ważna tylko dla tego jednego tekstu. Powtórz pomiar na innym, trudniejszym tekście, a zobaczysz różnicę. Lepiej jest mieć dwa pomiary niż jeden.
\end{przyklad}

Sam wynik nie wystarcza, bo można go wyśrubować, czytając za szybko. Dlatego do testu trzeba dołączyć sprawdzenie rozumienia. Zaraz po pomiarze, bez wracania do tekstu, zapisz trzy pytania o treść tych pięciu stron i odpowiedz na nie. Jeśli odpowiesz poprawnie na mniej niż dwa z trzech, nie zmierzyłeś tempa czytania, tylko tempo przerzucania słów. Zapisz wynik obok pomiaru: 240~sł./min, 2 z 3. Takie dwie liczby razem są prawdziwą miarą.

\begin{zasada}[Mierz na tekście, który czytasz zwykle]
Tempo zmienia się z gatunkiem. Ten sam czytelnik przy powieści, przy podręczniku z matematyki i przy artykule o prawie czyta z różnych prędkości. Jeżeli porównujesz wyniki, porównuj te same rodzaje tekstu. Jeśli zmierzysz się tylko raz, na lekkiej lekturze, zobaczysz liczbę, która opisuje lekturę, a nie czytanie.
\end{zasada}

\section{Co jest typowe, a co obiecują reklamy}

Dla dorosłego czytelnika tekstów nieliterackich, bez treningu, typowe tempo mieści się zwykle w zakresie od 200 do 300 słów na minutę\verify{typowe tempo czytania cichego dorosłych w tekstach nieliterackich mieści się zwykle w przedziale 200--300 słów na minutę; wartość zależy od tekstu i metody liczenia słów}. Ten przedział jest w tej książce przyjęty jako punkt odniesienia, a nie prawda o każdym czytelniku. Tabela w dalszej części rozdziału używa czterech wartości z tego zakresu: 150, 200, 250 i 300.

Przy 250 słowach na minutę typowa książka, licząca 80\,000 słów, zajmuje około \val{typical.hours} godziny czystego czytania. To liczba do zapamiętania: jeśli wiesz, ile godzin ma dana książka, wiesz, ile dni zajmie ci przy danym tempie i danym czasie.

Kursy szybkiego czytania obiecują zwykle wyraźnie więcej. Obietnice powyżej około 500 słów na minutę pojawiają się często, a towarzyszą im zwykle dwie rzeczy: pokaz z tekstem, który łatwo streścić, i brak niezależnego sprawdzenia rozumienia. Badania nad tym tematem, o ile ich wyniki da się w ogóle zebrać jednym zdaniem, nie potwierdzają, że tempo przekraczające kilkaset słów na minutę idzie w parze z rozumieniem porównywalnym z czytaniem wolniejszym\verify{przegląd badań nad szybkim czytaniem sugeruje, że wysokie tempo, zwykle powyżej około 500 słów na minutę, wiąże się z wyraźnie niższym rozumieniem; nie sprawdzono tego w źródłach pierwotnych dla tej książki}. Nie traktuj tego jako wyroku. Traktuj jako zasadę ostrożności: jeśli ktoś obiecuje tysiąc słów na minutę, zapytaj o test rozumienia, który sam przeprowadził, i zobacz, co zostaje po tygodniu.

Uczciwie: można nauczyć się przyspieszać w granicach tego, co już potrafisz. Przy lekturze prostej, znanej, z dobrą strukturą, wielu czytelników może czytać szybciej niż zwykle, bez utraty rozumienia. Granica przebiega tam, gdzie tekst wymaga myślenia. Tam przyspieszenie zwykle nie kupuje nic, bo to myślenie jest czasem, którego nie da się przeskoczyć.

\section{Co realnie daje czas}

Jeśli tempo jest w zakresie typowym, a wynik czytania i tak wydaje się za mały, poprawa przyjdzie z czterech rzeczy, które nie wymagają żadnej techniki czytania.

Pierwsza to \emph{czas}. Każda minuta dziennie, którą masz na czytanie, mnoży wynik przez tempo. Zanim zaczniesz martwić się o sekundy na stronie, sprawdź, czy masz te dziesięć minut. Rozdział~3 poświęcony jest temu, skąd je wziąć.

Druga to \emph{mniej przerwań}. Książka czytana w kawałkach, po dwie strony przy każdej okazji, kosztuje więcej, niż pokazuje licznik słów. Po każdej przerwie trzeba wrócić do miejsca, w którym się przerwało, odtworzyć wątek i przypomnieć sobie postacie albo argument. Ten koszt nie jest widoczny w pomiarze tempa, bo pomiar mierzy czas czytania, a nie czas odzyskiwania kontekstu. Jeśli czytasz w trzech sesjach po dziesięć minut, a nie w jednej sesji o trzydziestu minutach, płacisz za wejście w tekst trzy razy.

Trzecia to \emph{dobra lektura}. Książka, która nie wciąga, zostaje odłożona. Lektura, którą czytasz z przymusu, wymusza nie tylko przerwy, ale i przewracanie tych samych stron. Wybór tekstu, który pasuje do twojego poziomu i ciekawości, da ci więcej niż jakakolwiek technika. Rozdział~4 mówi o tym dokładniej.

Czwarta to \emph{cicha przestrzeń}. Hałas nie zmienia tego, ile słów mieści się w minucie, ale zmienia to, jak często uwaga ucieka od tekstu. Ciche miejsce nie musi być pokojem z drzwiami. Może być ławką w parku, krzesłem w kuchni po zmywaniu albo wyłączonym telefonem w innym pokoju. Warunek jest jeden: w czasie czytania nic nie wymaga od ciebie uwagi poza tekstem.

\section{Subwokalizacja bez mitów}

Subwokalizacja to wewnętrzne wymawianie słów w trakcie czytania. Wiele osób słyszy w głowie każde słowo, jakby ktoś je czytał na głos. Mit mówi, że trzeba się tego natychmiast pozbyć, bo to hamuje tempo do tempa mówienia. Jak wszystkie mity, ma w sobie ziarno prawdy i sporo nadużyć.

Ziarno jest takie: prawdopodobnie nie jesteś w stanie czytać znacznie szybciej niż wolno mówisz, jeśli każde słowo wymawiasz wewnątrz. Mówienie w tempie spokojnej rozmowy idzie wolniej niż zakres typowego czytania, więc przy tym mechanizmie sufit leży tam, gdzie leży tempo mowy\verify{typowe tempo mowy w spokojnej rozmowie to ok. 120--160 słów na minutę; zakres w języku polskim wymaga pomiaru, nie zakładamy go tu bez źródła}.

Nadużycie polega na tym, że ktoś wyciąga z tego wniosek, iż subwokalizacja jest zawsze błędem. Dla tekstów wymagających namysłu często nie jest, bo pozwala trzymać kolejne zdania w pamięci roboczej i sprawdzać, czy argument się trzyma. Wielu dobrych czytelników subwokalizuje świadomie i korzysta z tego. Czy jej brak poprawia rozumienie, zależy od tekstu i od czytelnika, a w tym wypadku mamy tylko częściowe dowody\verify{wyniki badań nad związkiem subwokalizacji z rozumieniem są niejednoznaczne; nie przytaczamy tu konkretnego eksperymentu}.

Praktyczna zasada jest więc taka: nie walcz z subwokalizacją w tekstach, które czytasz dla przyjemności, bo tam jej koszt jest mały. Jeśli chcesz ją ograniczyć, rób to świadomie, na krótkim fragmencie, i sprawdź rozumienie testem z poprzedniej części. Jeżeli wynik rozumienia spada, wróć do poprzedniego sposobu. Subwokalizacja, która nie przeszkadza, nie jest problemem, którego trzeba się pozbywać.

\begin{falle}[Ruch oczu jako miara umiejętności]
Część szkoleń uczy prowadzenia palcem lub wzrokiem po linii szybciej niż normalny ruch oczu. Efekt jest widoczny, bo oczy rzeczywiście przesuwają się szybciej. Pytanie brzmi, czy wzrok nadąża za informacją. Jeżeli przeskakujesz słowa, a potem nie pamiętasz zdania, nie przyspieszyłeś, tylko przestałeś czytać. Miarą jest test z pytaniami, nie szybkość przesuwania palca.
\end{falle}

\section{Trudne teksty czyta się wolno}

Istnieje rodzaj tekstu, przy którym każde wzmożenie tempa jest błędem. To teksty, w których każde zdanie opiera się na poprzednich: dowód, instrukcja, dłuższy wywód filozoficzny, rozdział o modelu statystycznym. Tam warto czytać powoli, z ołówkiem, zatrzymując się w miejscach, których nie rozumiesz.

Metoda ma trzy kroki. Najpierw przeczytaj cały fragment bez zatrzymywania się i bez zapisywania, żeby zobaczyć, dokąd zmierza. Potem przeczytaj go drugi raz, wolno, i przy każdym akapicie zapisz na marginesie jedno zdanie własnymi słowami, co on twierdzi. Na koniec sprawdź, czy potrafisz wyjaśnić główną myśl rozdziału bez zaglądania do tekstu. Jeśli nie potrafisz, znalazłeś miejsce, które trzeba przeczytać jeszcze raz, a nie miejsce, w którym czytanie się nie udało.

Ołówek nie jest ozdobą. Zmusza do wyrażenia myśli zamiast jej tylko rozpoznania, a to różnica, którą pokażą pytania po tekście. Tekst przeczytany ze zrozumieniem zostawia ślad, który da się opisać. Tekst przeczytany szybko zostawia wrażenie.

Tempo w takich lekturach liczy się inaczej. Pięć stron trudnego tekstu może zająć dwadzieścia minut i nie oznacza to porażki. Jeśli licznik pokazuje 80 słów na minutę przy wykresie i wzorze, to jest wynik właściwy dla tego tekstu. Zanotuj go, ale nie porównuj z wynikiem z lekkiej prozy.

\begin{uwaga}[Kiedy zmienić tempo]
Sygnałem do zwolnienia nie jest uczucie, że tekst jest trudny, tylko to, że ostatni akapit wydaje się znany, a nie pamiętasz z niego niczego. Wtedy czytasz bez czytania. Zatrzymaj się, cofnij o akapit i przeczytaj wolniej. To jest zmiana tempa, którą warto mieć pod ręką.
\end{uwaga}

\section{Tryb skupienia i rachunek}

Jedną z prostych rzeczy, które realnie pomagają, jest tryb skupienia: blok czytania o stałej długości, bez telefonu w zasięgu ręki i z jednym tekstem. Wzorem jest blok dwudziestopięciominutowy, po którym następuje krótka przerwa. Protokół jest prosty: ustaw stoper na 25 minut, odłóż telefon do innego pokoju, czytaj jedną książkę, a po zakończeniu bloku zapisz numer strony i jedno zdanie o tym, co z niego zostało. Przerwa to pięć minut, w których nie wracasz do ekranu, tylko wstajesz.

Co do efektów, nie obiecujemy cudów. Trzeba uczciwie zaznaczyć, że wpływ takich bloków na długość czytania w skali roku nie został tu zmierzony i jest zwykle przedmiotem badań o niejednoznacznych wynikach\verify{efekty krótkich bloków skupienia na długość i jakość czytania nie zostały zmierzone w tej książce; wpływ metod zarządzania uwagą jest przedmiotem badań o niejednoznacznych wynikach}. Protokół warto wypróbować jako eksperyment osobisty: przez dwa tygodnie policz, ile stron przeczytałeś w blokach, i porównaj z poprzednimi dwoma tygodniami.

Teraz rachunek, który łączy wszystko. Tabela pokazuje liczbę książek w roku dla różnych kombinacji minut dziennie i tempa.

\generated{books-per-year}

Weźmy trzydzieści minut dziennie przy 200 słowach na minutę. Wychodzi \val{books.30.200} książki rocznie. Co daje lepszy wynik: wzrost tempa o 50 słów na minutę czy dodatkowe dziesięć minut dziennie?

\begin{przyklad}[Dwie dźwignie]
Wzrost tempa z 200 do 250 słów na minutę przy 30 minutach dziennie:
$50 \cdot 30 \cdot 365 = 547\,500$ słów rocznie, czyli $547\,500 / 80\,000 \approx 6{,}8$ książki. Z tabeli: \val{books.30.250} minus \val{books.30.200}.

Dziesięć minut więcej przy 200 słowach na minutę:
$10 \cdot 200 \cdot 365 = 730\,000$ słów rocznie, czyli $730\,000 / 80\,000 \approx 9{,}1$ książki.

Przy tym punkcie wyjścia dodatkowe dziesięć minut daje więcej niż pięćdziesiąt słów więcej na minutę. Ale zauważ, od czego zależy wynik. Oba wzrosty są równie cenne, gdy $50 \cdot M = 10 \cdot W$, czyli gdy $W = 5M$. Przy $M = 30$ minutach dzieje się to dla $W = 150$ słów. Sprawdźmy: dziesięć dodatkowych minut przy 150 słowach daje $\val{books.30.150}$ minus $\val{books.20.150}$, czyli około 6,8 książki, a wzrost z 150 do 200 słów przy 30 minutach daje $\val{books.30.200}$ minus $\val{books.30.150}$, czyli również około 6,8. Przy wolniejszym czytelniku czas wygrywa z tempem, przy szybszym tempo zaczyna wygrywać z czasem.
\end{przyklad}

Wniosek praktyczny jest prosty. Dopóki nie masz czasu, poprawa tempa daje mało, bo nie ma czego mnożyć. Dopóki tempo jest typowe, dziesięć minut dziennie jest tańszą dźwignią niż trening czytania. Dopiero gdy masz czas i czytasz w zakresie typowym, a chcesz więcej, warto sięgnąć po tempo, i to w granicach, które sprawdziłeś testem z rozumieniem.

Dla porządku: w tym rachunku książka ma 80\,000 słów. Dla krótszej książki czas zakupuje więcej tytułów, dla dłuższej mniej. Na tym tle liczba \val{need.24}, czyli minut dziennie na dwadzieścia cztery książki rocznie przy 250 słowach na minutę, pokazuje, jak mało czasu wystarcza, jeśli plan jest konkretny.

\begin{cwiczenia}
  \exitem{Oblicz tempo z danych: 1\,200 słów przeczytanych w 5 minut 0 sekund.}{$1\,200 / 5 = 240$ słów na minutę. Zapisz obok wynik testu rozumienia, bo sam wynik tempa nie wystarcza.}
  \exitem{Ile słów przeczytasz w ciągu 20 minut przy 230 słowach na minutę?}{$230 \cdot 20 = 4\,600$ słów.}
  \exitem{Ile książek rocznie zmieści się przy 30 minutach dziennie i 250 słowach na minutę? Wynik podaj z dokładnością do jednego miejsca po przecinku.}{$30 \cdot 250 \cdot 365 = 2\,737\,500$ słów; $2\,737\,500 / 80\,000 \approx 34{,}2$ książki. Zgadza się z tabelą: \val{books.30.250}.}
  \exitem{Dlaczego dziesięć dodatkowych minut przy 150 słowach na minutę daje tyle samo, co wzrost tempa z 150 do 200 słów na minutę przy 30 minutach?}{Oba przypadki dają po $1\,500$ słów dziennie: $10 \cdot 150 = 1\,500$ oraz $50 \cdot 30 = 1\,500$. Dlatego w tabeli różnica między 20 a 30 minutami przy 150 słowach jest równa różnicy między 150 a 200 słowach przy 30 minutach.}
  \exitem{Po pięciu stronach trudnego eseju zadajesz sobie trzy pytania i odpowiadasz na jedno. Co z tego wynika dla pomiaru?}{Wynik rozumienia to 1 z 3, więc nie należy zapisywać tempa jako miary czytania. Trzeba zwolnić i powtórzyć pomiar na tym samym tekście, z ołówkiem. Tempo bez rozumienia nie jest miarą sprawności.}
  \exitem{Kursy obiecują tysiąc słów na minutę. Co powinieneś zażądać, zanim w to uwierzysz?}{Niezależnego testu rozumienia na tekście, którego nie znałeś wcześniej, z pytaniami przygotowanymi przed czytaniem i wynikiem podanym razem z tempem. Bez tego obietnica jest tylko liczbą.}
  \exitem{Ile minut dziennie potrzebujesz na 24 książki rocznie przy 250 słowach na minutę?}{$24 \cdot 80\,000 / 250 / 365 \approx 21{,}0$ minuty dziennie. W tabeli to wartość \val{need.24}.}
\end{cwiczenia}

\zadanie{Zmierz tempo na pięciu pełnych stronach książki, którą czytasz teraz. Zapisz liczbę słów, czas w sekundach, wynik w słowach na minutę i wynik trzech pytań o treść. Powtórz pomiar trzy razy w ciągu tygodnia, w tym raz na tekście trudnym. Sukces to trzy zapisane pomiary z trzema wynikami rozumienia, a także jeden blok skupienia dwudziestopięciominutowy dziennie przez pięć dni z zapisaną liczbą przeczytanych stron.}
